\documentclass[../../main.tex]{subfiles}

\begin{document}

\section{Subspaces}

In the previous section, we discussed that all ten axioms must be
satisfied for a set to be a vector space. However, testing and
proving that a set satisfies all ten axioms, to be honest, is quite
tedious. The good news is that we only need to prove three conditions
to show that a set is a vector space thanks to \textbf{subspaces}!

\begin{definition}{}{}
A \textbf{subspace} of a vector space $V$ is a subset of $V$ that uses
the same definition of vector addition and scalar multiplication from
$V$.
\end{definition}

Now for a set to be a subspace, it must satisfy the following
conditions.

\begin{theorem}{}{Conditions of a Subspace}
A non-empty subset $S$ of a vector space $V$ is a subspace of $V$ if
and only if the following conditions are satisfied.
\begin{enumerate}[noitemsep]
\item
\textbf{Closure for Addition: }
$\forall u, v \in S, u + v \in S$
\item
\textbf{Closure for Scalar Multiplication: }
$\forall u \in S$ and scalars $\alpha$, $\alpha u \in S$
\end{enumerate}
\end{theorem}
\begin{proof}
Notice that it suffices to show that $S$ is a vector space, or adhere
to the then axioms for a vector space. Moreover, A1 and S1 are assumed
to be true by the initial condition.

By the second condition, $\mathbf{0} = 0 u \in S$ and $-u = (-1) u
\in S$. Therefore, A4 and A5 are satisfied. For A2 and A3, $u, v, w
\in V$ is true for $u, v, w \in S$ by definition. Moreover, because
$V$ is a vector space, A2 and A3 hold. Similarly, because $V$ is a
vector space that satisfies S1, S2, S3, S4, and S5, its subset $S$
also satisfies the conditions. Therefore, if the two closure
conditions are met and $S$ is a subset of $V$, then $S$ is a subspace
of $V$.
\end{proof}

Now the good thing is that subspaces are vector spaces, so it suffices
to show that a set is a subspace of a known vector space with the two
conditions shown above to show that it is a vector space. Let's take
a look at an example.

\begin{problem}{}{}
Determine whether $S = \{ (x, y, z) \in \mathbb{R}^3 \mid y = 0 \}$
is a vector space.
\end{problem}
\begin{solution}
Notice that $S$ is a non-empty subset of $\mathbb{R}^3$, which is a
known vector space. Therefore, it suffices to show that $u =
(x_1, 0, z_1), v = (x_2, 0, z_2) \in S$ and scalars $\alpha$ satisfy
the closure conditions for addition and scalar multiplication.

Notice that $u + v = (x_1 + x_2, 0, z_1 + z_2) \in S$. Moreover,
$\alpha u = (\alpha x_1, 0, \alpha z_1) \in S$. Therefore, $S$ is a
vector space, more specifically a subset of $\mathbb{R}^3$.
\end{solution}

As you can see from the solution above, subspaces really made our
lives easier. However, to make it even better, we can boil it down
to a single condition.

\begin{corollary}{}{The Condition for a Subspace}
For scalars $\alpha$ and $\beta$ and $u, v \in S$ where $S$ is a
non-empty subset of a vector space $V$ that has the same operations
for vector addition and scalar multiplication, $S$ is a subspace if
and only if the following condition is satisfied.
\[
\alpha u + \beta v \in S
\]
\end{corollary}
\begin{proof}
First and foremost, the first half of the statement can be proven.
If $\alpha u + \beta v \in S$, then $u + v \in S$ where $\alpha =
\beta = 1$ since $\alpha$ and $\beta$ are all scalars. Similarly,
$\alpha u + \beta v \in S$ implies that $\alpha u \in S$ where
$\beta = 0$. By Theorem \ref{theo:Conditions of a Subspace}, $S$ is a
subspace of $V$.

Continuing with the second half of the statement, if $S$ is a subset
of $V$, then $\alpha u, \beta v \in S$. Therefore, $\alpha u +
\beta v \in S$, proving the corollary.
\end{proof}

More examples on this can be found in the practice problems below.
Now let's get into the topics of \textbf{linear combination} and the
\textbf{span} of a set.

\subsection{Linear Combinations and Span}

As always, let's start with definitions.

\begin{definition}{}{}
For scalars $\alpha_i$ and $v_i$ in a vector space $V$, the vector
$u$ defined as the following is a \textbf{linear combination of the
vectors} $v_i$.
\[
u = \sum_{k = 1}^n \alpha_k v_k
\]
\end{definition}

For instance, $[1, 2, 3]$ is a linear combination of $[2, -2, 1]$,
$[1, 0, 2]$, and $[1, -1, 1]$ as it can be represented as the
following.
\[
\begin{bmatrix}
    1 & 2 & 3
\end{bmatrix}
=
\begin{bmatrix}
    2 & -2 & 1
\end{bmatrix}
+
3\begin{bmatrix}
    1 & 0 & 2
\end{bmatrix}
-
4\begin{bmatrix}
    1 & -1 & 1
\end{bmatrix}
\]
Now we have a special term for the collection of such vectors.

\begin{definition}{}{}
The \textbf{span of the set of vectors} $\{ v_1, \ldots v_n \}$ is
the set of all possible linear combinations of the vectors.
\[
\Span\{v_1, \ldots, v_n\}
= \left\{
    \sum_{k = 1}^n \alpha_k v_k \; \middle|\; \alpha_k \in \mathbb{R}
\right\}
\]
\end{definition}

With the definition, we could notice an important theorem.

\begin{theorem}{}{Span and Subspace I}
For a vector space $V$ and its subset $S = \{ v_1, \ldots, v_n \}$,
$\Span(S)$ is a subspace of $V$.
\end{theorem}
\begin{proof}
First and foremost, notice that $\mathbf{0} \in \Span(S)$ since all
the scalars in linear combination can equal zero. Moreover, let $u =
\sum_{k = 1}^n \alpha_k v_k$ and $v = \sum_{k = 1}^n \alpha_k v_k$ for
any scalars $\alpha_k$ and $\beta_k$. By definition, $\alpha u +
\beta v \in \Span(S)$ is true. Therefore, by Corollary
\ref{cor:The Condition for a Subspace}, $\Span(S)$ is a subspace of
$V$.
\end{proof}

Side note here is that there are spans of infinite sets and they
adhere to the theorem above. Continuing with the definition of the
span of a set and its relation with subspace, we can show the two
following theorems.

\begin{theorem}{}{Span and Subspace II}
Consider the following statements for a vector space $V$ and its
subspace $W$ and non-empty subset $S$ such that $S \subset W$.
\begin{enumerate}[noitemsep]
\item $\Span(S) \subset W$
\item The intersection of all subspaces containing $S$ is $\Span(S)$.
\end{enumerate}
\end{theorem}

Let's start by proving the first statement.

\begin{proof}
Let $u = \sum_{k = 1}^n \alpha_k v_k \in \Span(S)$. By definition,
$v_k \in W$ for integer $k \in [1, n]$. Moreover, $u \in W$ by
Corollary \ref{cor:The Condition for a Subspace} and $\Span(S)
\subset W$.
\end{proof}

Below is the proof for the second one.

\begin{proof}
First and foremost, notice that by the previous statement, it was
shown that $\Span(S)$ is included in all subspaces that contains $S$.
Let $A$ be the intersection of all subspaces that contain $S$.
Therefore, $\Span(S) \subset A$. Moreover, by Theorem
\ref{theo:Span and Subspace I} $\Span(S)$ is a subspace of $V$ that
contains $S$. Thus, $A \subset \Span(S)$. Because $\Span(S) \subset A$
and $A \subset \Span(S)$, $\Span(S) = A$.
\end{proof}

Next, let's conclude our discussion on linear combination and span
of a set of vectors by proving the following statements.

\begin{theorem}{}{Four Lemmas on Span of a Set of Vectors}
For a vector space $V$ and its subsets $S, T \subset V$, consider the
following statements.
\begin{enumerate}[noitemsep]
\item If $T \subset S$, then $\Span(T) \subset \Span(S)$.
\item $\Span(\Span(S)) = \Span(S)$.
\item If $T \subset \Span(S)$, then $\Span(T) \subset \Span(S)$.
\item For nonzero vectors $v_1, \ldots, v_n \in V$, scalar
$\lambda \in \mathbb{R}$, and $1 \leq j, k \leq n$ such that $\lambda
\neq 0$ and $j \neq k$, the following equations hold.
\vspace{-0.5\baselineskip}
\begin{align*}
    \Span\{ v_1, \ldots, v_k, \ldots, v_n \}
    &= \Span\{ v_1, \ldots, v_k + \lambda v_j, \ldots, v_n \} \\
    \Span\{ v_1, \ldots, v_k, \ldots, v_n \}
    &= \Span\{ v_1, \ldots, \lambda v_k, \ldots, v_n \}
\end{align*}
\end{enumerate}
\end{theorem}

There's a lot to prove, but let's do them one by one starting with
the first one.

\begin{proof}
Notice that since all $t_k \in T$ satisfy $t_k \in S$, all
$\sum_{k=1}^n c_k t_k \in \Span(T)$ for some appropriate $n$ and
constant $c_k$ also satisfy $\sum_{k=1}^n c_k t_k \in \Span(S)$.
Therefore, $\Span(T) \subset \Span(S)$.
\end{proof}

Below is the proof for the second statement, which is a corollary
of the first statement.

\begin{proof}
First, notice that $S \subset \Span(S)$. By the first statement,
$\Span(S) \subset \Span(\Span(S))$ holds. Moreover, notice that
because $\Span(S)$ contains all linear combinations of $S$,
making more linear combinations of the linear combinations of $S$
will not add more elements. Therefore, $\Span(\Span(S)) \subset
\Span(S)$ and $\Span(\Span(S)) = \Span(S)$.
\end{proof}

An interesting fact about this property is that this applies for
any subset $S$, finite or infinite. This is because even if $\Span(S)$
is infinite, it only uses finite number of vectors to create
each linear combination. Continuing, let's prove the third statement.

\begin{proof}
Notice that by the first statement, $\Span(T) \subset
\Span(\Span(S))$. Continuing with the second statement shows that
$\Span(T) \subset \Span(S)$.
\end{proof}

Finally, here is the proof for the last statement.

\begin{proof}
First, it is self-evident that $\Span\{ v_1, \ldots, v_k
+ \lambda v_j, \ldots, v_n \} \subset \Span\{ v_1, \ldots,
v_k, \ldots, v_n \}$ since all linear combination of the set of
vectors on the left can be represented as the linear combination of
the set of vectors on the right. Moreover notice that $v_k$ can be
represented as $v_k = (v_k + \lambda v_j) - \lambda v_j$ even if
$k = j$. Therefore, $v_k \in \Span\{ v_1, \ldots, v_k + \lambda v_j,
\ldots, v_n \}$ and $\Span\{ v_1, \ldots, v_k, \ldots, v_n \} \subset
\Span\{ v_1, \ldots, v_k + \lambda v_j, \ldots, v_n \}$ by
statement 3. Therefore, the equations in the statement hold.
\end{proof}

With the ideas of linear combinations and the span of a set of vectors
in mind, let's discuss about linear independence.

\subsection{Linear Independence}

One of the most important reasons why we use Linear Algebra is due to
efficiency. As you could have guessed, most vector spaces have
infinitely many vectors. Of course, there must be an efficient way
of representing the vector space instead of manually listing.
Determining such representative vectors involves numerous factors,
and that's what we will discuss with linear independence, basis,
and dimension.

\begin{definition}{}{}
Consider the linear combination of $V = \{ v_1, \ldots, v_n \}$
in a vector space and scalars $c_1, \ldots, c_n$ that satisfy the
following equation.
\[
\sum_{k = 1}^n c_k v_k = \mathbf{0}
\]
The set of vectors $V$ is \textbf{linearly dependent} if there exists
scalars $c_1, \ldots, c_k$ that are not all equal to zero. The set of
vectors is \textbf{linearly independent} if all scalars are zero in
the equation.
\end{definition}

Let's take a look at a quick example.

\begin{problem}{}{}
Determine whether $V = \{ \langle 1, 1 \rangle, \langle 2, 3\rangle
\in \mathbb{R}^2 \}$ is linearly independent.
\end{problem}
\begin{solution}
By definition, for the set to be linearly independent, all scalars
$c_1$ and $c_2$ must be zero for the following equation to hold.
\[
c_1 \langle 1, 1 \rangle + c_2 \langle 2, 3 \rangle
= \langle 0, 0 \rangle
\]
Writing in system of linear equations, the equation above can
be written as the following.
\begin{align*}
    c_1 + 2c_2 &= 0 \\
    c_1 + 3c_2 &= 0
\end{align*}
Solving the equation, $c_1 = c_2 = 0$ and no other set of scalars
can lead to the equation. Therefore, the set of vectors is linearly
independent.
\end{solution}

One interesting observation that we can make is that for $n \in
\mathbb{Z}$, if the set of vectors in $\mathbb{R}^n$ contains more
than $n$ elements, then the set is linearly dependent. We will
discuss more on this later when we discuss matrix and its application
in linear systems, but for such cases, we have more variables than
linear equations, which implies that there exists infinitely many
solutions. Therefore, such a set of vectors must be linearly
dependent. Continuing, below is a theorem that we can deduce from
the definition.

\begin{theorem}{}{Finite Set of Vectors and Linear Dependence}
A set with a finite number of nonzero vectors is linearly dependent
if and only if a vector in the ordering of the set can be represented
as a linear combination of vectors preceding it.
\end{theorem}
\begin{proof}
First and foremost, the first half of the theorem can be shown.
Consider the set of vectors $\{ v_1, \ldots, v_{j-1}, v_j, \ldots,
v_n \}$. Assume for scalars $c_k$, the following equation is true.
\[
v_j
= \sum_{k = 1}^{j-1} c_k v_k
\]
Therefore, $\mathbf{0} = \sum_{k = 1}^{j-1} c_k v_k - v_j$.
Letting $c_k = 0$ for integer $k \in (j, n]$, the following equation
holds.
\[
\sum_{k = 1}^{j-1} c_k v_k - v_j + \sum_{k = j+1}^{n} c_k v_k
= \mathbf{0}
\]
Because at least one scalar is nonzero, the set of vectors $\{ v_1,
\ldots, v_n \}$ is linearly dependent.

Continuing, let the set of vectors $\{ v_1, \ldots, v_n \}$ be
linearly dependent. For the corresponding scalars $c_k$, let
$c_j$ be the last scalar that is nonzero. Thus, the following
equation holds.
\[
\sum_{k = 1}^{j} c_k v_k + \sum_{k = j+1}^{n} 0 v_k
= \mathbf{0}
\]
Rearranging the equation, $v_k = -\sum_{k = 1}^{j-1} \frac{c_k}{v_j}
v_k$, which is a linear combination. Because the converse holds,
the theorem is true.
\end{proof}

From this theorem, we can deduce a corollary.

\begin{corollary}{}{}
For a vector space $V = \Span(S)$ and linearly dependent set
$S = \{ v_1, \ldots, v_n \} \subset V$, then there exists a set
$T \subset S$ such that $|T| = n - 1$ and $V = \Span(T)$.
\end{corollary}
\begin{proof}
By Theorem \ref{theo:Finite Set of Vectors and Linear Dependence},
there exists element $v_j \in S$ such that $v_j$ is the linear
combination of the previous elements as listed in $S$.

Let $T = \{ v_1, \ldots, v_{j-1}, v_{j+1}, \ldots, v_n \}$.
Consider the following equations for scalars $d_k$ and $c_k$.
\begin{align*}
    V
    &= \sum_{k = 1}^n d_k v_k
    = \sum_{k = 1}^{j-1} d_k v_k + d_j v_j +
        \sum_{k = j+1}^n d_k v_k \\
    &= \sum_{k = 1}^{j-1} d_k v_k +
        d_j \left( \sum_{k = 1}^{j-1} c_k v_k \right) +
        \sum_{k = j+1}^n d_k v_k \\
    &= \sum_{k = 1}^{j-1} (d_j c_k + d_k) v_k +
        \sum_{k = j+1}^n d_k v_k
    = \Span(T)
\end{align*}
Therefore, the corollary holds.
\end{proof}

Continuing from the corollary, the following theorems could be
asserted.

\begin{theorem}{}{Properties of Linear Independence and Dependence}
Consider the following statements for set $S$.
\begin{enumerate}[noitemsep]
\item
If $S$ contains a zero vector, then it is linearly dependent.
\item
If $|S| = 1$, then it is linearly dependent if and only if its
element is the zero vector.
\item
If $|S| = 2$, then it is linearly dependent if and only if a vector
is a scalar multiple of the other.
\item
If $S$ is linearly dependent, then every set that contains $S$ is
linearly dependent.
\item
If $S$ is linearly independent, then all its subsets are linearly
independent.
\end{enumerate}
\end{theorem}

There's quite a bit to prove! Let's start with the first one.

\begin{proof}
Let $S = \{ \mathbf{0}, v_1, \ldots, v_n \}$. Notice that
$\alpha(\mathbf{0}) + \sum_{k = 1}^n 0 v_k = \mathbf{0}$ for any
scalar $\alpha$. Therefore, $S$ is linearly dependent.
\end{proof}

Below is the proof for the second one.

\begin{proof}
Let $S = \{ v \}$. Notice that if $v \neq \mathbf{0}$, it is not
possible for a nonzero scalar $\alpha$ to make $\alpha v =
\mathbf{0}$. Therefore, if $|S| = 1$, then $v = \mathbf{0}$.
For the converse, if $v = \mathbf{0}$, then any scalar $\alpha$
will let $\alpha v = \mathbf{0}$. Therefore the statement holds.
\end{proof}

Here is the proof for the third statement.

\begin{proof}
Let $S = \{ v_1, v_2 \}$. First, if $v_2 = \alpha v_1$ for any scalar
$\alpha$, then $\mathbf{0} = \alpha v_1 - \alpha v_1 = \alpha v_1 -
v_2$. Therefore, $S$ is linearly dependent. For the converse,
$c_1 v_1 + c_2 v_2 = 0$ for scalars $c_1$ and $c_2$. Therefore,
assuming that the scalars are nonzero without loss of generality,
$v_2 = -\frac{c_1}{c_2} v_1$. Thus, the statement is true.
\end{proof}

Below is the proof for the fourth one.

\begin{proof}
Let $S = \{ v_1, \ldots, v_n \}$ and $T = \{ u_1, \ldots, u_n,
v_1, \ldots, v_n \}$. By definition, there exists $c_k$ that are
not all zero such that $\sum_{k = 1}^n c_k v_k = \mathbf{0}$.
Because $\sum_{k = 1}^n 0 u_k + \sum_{k = 1}^n c_k v_k = \mathbf{0}$,
$T$ is also linearly dependent.
\end{proof}

Finally, here is the proof for the last statement.

\begin{proof}
Let $T \subset S$ where $S$ is linearly independent. Assume for
the sake of contradiction that there exists a subset $T$ that is
linearly dependent. However, if $T$ is linearly dependent, then
$S$ is linearly dependent by the previous property, which is
a contradiction. Therefore, by contradiction, if $S$ is linearly
independent then all its subsets $T$ are also linearly independent.
\end{proof}

With this, we have concluded our discussion on linear dependence
and independence. Before concluding this note, let's discuss
basis and dimension of a set of vectors.

\subsection{Basis and Dimension}

As always, let's start with definition.

\begin{definition}{}{}
For a vector space $V$ and its subset $S$, if $\Span(S) = V$, then
$S$ is a spanning set for $V$.
\end{definition}

When we say aloud, we say that $S$ spans $V$. For instance, the set
$\{ \langle 1, 0 \rangle, \langle 0, 1 \rangle \}$ spans
$\mathbb{R}^2$ since all vectors $\langle a, b \rangle \in
\mathbb{R}^2$ can be represented as the sum $a \langle 1, 0 \rangle
+ b\langle 0, 1 \rangle$.

\begin{definition}{}{}
A linearly independent subset that spans $V$ is a \textbf{basis} for
$V$.
\end{definition}

Continuing from the previous example, we can see that $\{ \langle 1, 0
\rangle, \langle 0, 1 \rangle \}$ not only spans $\mathbb{R}^2$,
but is also a basis for $\mathbb{R}^2$ since it is linearly
independent. As you can see, the same will apply if we generalize it
to $\mathbb{R}^n$.

\begin{definition}{}{}
The \textbf{standard basis for $\mathbb{R}^n$} is the set $S =
\{ e_1, \ldots e_n \}$ where $e_i$ is an element of $\mathbb{R}^n$
with $1$ in its $i^\text{th}$ entry and $0$ elsewhere.
\end{definition}

We could do the same for polynomials. First, notice that the set
$\{ 1, x, x^2 \}$ spans $\mathbb{P}^2$. Moreover, the set is linearly
independent. Therefore, $\{ 1, x, x^2 \}$ is a basis for
$\mathbb{P}^2$.

\begin{definition}{}{}
The \textbf{standard basis for $\mathbb{P}^n$} is the set $S =
\{ 1, x, x^2, \ldots, x^n \}$.
\end{definition}

From the two definitions above, we can notice that there are finite
elements for the two standard bases for $\mathbb{R}^n$ and
$\mathbb{P}^n$. We call such vector spaces
\textbf{finite-dimensional}.

\begin{definition}{}{}
A vector space with a basis of finite elements is called
\textbf{finite-dimensional}. Whereas, \textbf{infinite dimensional}
vector spaces do not contain a basis with finite elements.
\end{definition}

An example of infinite dimensional vector space would be $\mathbb{P}$
since the basis is required to be an infinite set. However, we will
not be discussing infinite dimensional vector spaces. From the
definitions above, we could establish a theorem.

\begin{theorem}{}{Basis and Linear Dependence}
For a basis $S$ of a vector space $V$ such that $|S| = n$, any set
in $V$ with more than $n$ elements is linearly dependent.
\end{theorem}
\begin{proof}
Let $S = \{ v_1, \ldots, v_n \}$ and $T = \{ u_1, \ldots, u_m \}$
such that $S, T \in V$ and $m > n$. Notice that because $\Span(S) = V$
and $T \in V$, there exists a set of constants $a_{ij}$ that satisfies
the following equation (note that we will be using this indices in
the next note when we discuss matrices.).
\begin{align*}
    u_1 &= \sum_{j=1}^n a_{1j} v_j \\
    u_2 &= \sum_{j=1}^n a_{2j} v_j \\
    &\quad\vdots \\
    u_m &= \sum_{j=1}^n a_{mj} v_j
\end{align*}
Therefore, it suffices to show that there exists constants $c_1,
\ldots, c_m$ such that not all are zero and that satisfies the
following equation.
\[
\sum_{i=1}^m \sum_{j=1}^n c_i a_{ij} v_j
= \mathbf{0}
\]
Continuing, it suffices to show that the variables $c_i$ and constants
$a_{ij}$ satisfy the following equations where not all $c_i$ are zero.
\begin{align*}
    \sum_{i=1}^m c_i a_{i1} &= 0 \\
    \sum_{i=1}^m c_i a_{i2} &= 0 \\
    &\ \vdots \\
    \sum_{i=1}^m c_i a_{in} &= 0
\end{align*}
Notice that there are $m$ variables and $n$ equations. It is known
that for a linear system of equations, there exists an infinite number
of solutions if there are more variables than the equations (don't
worry we will prove this in future notes.). Because $m > n$, there
exists $c_1, \ldots c_m$ that are not all zero and that satisfy the
equations for linear dependence, the theorem holds.
\end{proof}

Continuing, we can derive three corollaries from the theorem.

\begin{corollary}{}{}
If a basis for the vector space $V$ contains $n$ elements, then any
linearly independent subsets of $V$ contain at most $n$ elements.
\end{corollary}
\begin{proof}
By definition, a basis of a vector space $V$ spans $V$. Because any
linearly dependent subsets of $V$ contain more than $n$ elements
by Theorem \ref{theo:Basis and Linear Dependence}, any linearly
independent subsets of $V$ must not contain more than $n$ elements.
\end{proof}

Below is the second corollary.

\begin{corollary}{}{}
If $V$ is a finite vector space, every basis of $V$ will contain
the same number of elements.
\end{corollary}
\begin{proof}
Let $S$ and $T$ be bases of $V$ with $p$ and $q$ elements
respectively. By definition, the bases of a vector space are linearly
independent. By the previous corollary, it is evident that $q \leq p$
and $p \leq q$ are true. Therefore, $p = q$ and the corollary holds.
\end{proof}

Now before continuing with our third corollary, let's first define
what the dimension of a vector space is.

\begin{definition}{}{}
For a finite dimensional vector space $V$, the number of elements
in any basis of $V$ is called the \textbf{dimension} of $V$,
denoted as $\dim(V)$.
\end{definition}

For instance, recall the standard basis of $\mathbb{R}^n$. From the
number of elements in the basis, we can write that $\dim(\mathbb{R}^n)
= n$. Similarly, we can also notice that $\dim(\mathbb{P}^n) = n + 1$.
Now let's continue with our third corollary.

\begin{corollary}{}{}
Every set $S$ in an $n$-dimensional vector space such that $|S| =
n + 1$ is linearly dependent.
\end{corollary}
\begin{proof}
By definition, a basis of an $n$-dimensional vector space $V$ contains
$n$ elements. By Theorem \ref{theo:Basis and Linear Dependence}, since
$n + 1 > n$, $S$ must be linearly dependent.
\end{proof}

One fun fact to note is that for the vector space $\{0\}$, we say
$\dim(\{0\}) = 0$ by convention. Moreover, consider the following
theorem.

\begin{theorem}{}{Unique Linear Combination}
For a basis $S = \{ v_1, \ldots, v_n \}$ of a vector space $V$,
any vector $u \in V$ can be written as a unique linear combination
of $S$.
\end{theorem}
\begin{proof}
For the sake of contradiction, assume that $u$ can be written as
$u = \sum_{k=1}^n c_k v_k = \sum_{k=1}^n d_k v_k$ where $c_k \neq d_k$
for some choice of $k$. Therefore, the following equations can be
obtained.
\begin{align*}
    \sum_{k=1}^n c_k v_k &= \sum_{k=1}^n d_k v_k \\
    \sum_{k=1}^n (c_k - d_k) v_k &= 0
\end{align*}
Note that by definition, $S$ is linearly independent and $c_k - d_k
= 0$ for all integer $k \in [1, n]$. Because it contradicts with
the assumption that $c_k \neq d_k$ for some choice of $k$, the theorem
holds.
\end{proof}

Now with the theorem in mind, because we know that the linear
combination of a basis to represent a vector is unique, we can
establish the following definition.

\begin{definition}{}{}
For a basis $S = \{ v_1, \ldots, v_n \}$ of a vector space $V$
and any vector $u = \sum_{k=1}^n c_k v_k \in V$, the
\textbf{coordinates of $u$ with respect to $S$}, denoted as $(u)_S$
is the list of coefficients $c_1, \ldots, c_n$.
\end{definition}

We often represent the coordinates with $n$-tuple. To see that,
let's take a look at the following example.

\begin{problem}{}{}
Find the coordinate of $f(x) = x^2 + 2x + 3$ with respect to the
basis $S = \{ x^2, 2x^2 + 1, x \}$ in the vector space $\mathbb{P}^2$.
\end{problem}
\begin{solution}
Consider the following equation.
\[
f(x)
= 3 \left( 2x^2 + 1 \right) - 5 \left( x^2 \right) + 2(x)
\]
Therefore, the coordinate of $f(x)$ can be represented as the
following $3$-tuple.
\[
\begin{bmatrix}
    -5 \\
    3 \\
    2
\end{bmatrix}_S
\]
\end{solution}

Lastly, let's conclude our discussion on basis and dimension with
a theorem.

\begin{theorem}{}{Properties of Subset}
Consider the following statements for a vector space $V$ and
its subset $S \subset V$.
\begin{enumerate}[noitemsep]
\item
There exists a basis for $V$ that is a subset of $S$ if $\Span(S) =
V$.
\item
There exists a basis for $V$ that includes all elements in $S$ if
$S$ is linearly independent.
\end{enumerate}
\end{theorem}

Let's start by proving the first statement.

\begin{proof}
Notice that by definition, a basis of a vector space $V$ is a linearly
independent set that spans $V$. Therefore, it suffices to show that
there exists a linearly independent set that spans $V$ that is also
a subset of $S$.

By Theorem \ref{theo:Finite Set of Vectors and Linear Dependence},
a subset of $S$ is linearly dependent if any only if a vector can be
represented as the linear combination of the preceding vectors.
Let $T$ be a subset of $S$ without the vectors that can be represented
as the linear combination of the preceding vectors. Notice that $T$
spans $V$ as the vector removed can be represented as the linear
combination of the other vectors. Moreover, because $T$ is linearly
independent, there exists a basis for $V$ that is a subset of $S$.
\end{proof}

Below is the proof for the second statement.

\begin{proof}
Let $T$ be a basis for $V$ and $A \coloneqq S \cup T$. By definition,
$\Span(A) = V$. Construct $B \in A$ by removing all vectors that can
be represented as a linear combination of the preceding vectors.
Notice that no elements from $S$ will be removed as $S$ is linearly
independent. By construction, $B$ spans $V$ and is also linearly
independent. Because $S \subset B$, the statement holds.
\end{proof}

We discussed fundamental topics of vector spaces. Of course there
are a lot more to discuss in vector spaces; however, to do so,
we need some understanding of matrices. So let's pause our discussions
on vector spaces for now, and move on to matrices and their
applications in linear system of equations before looking at their
implications in vector spaces.

\end{document}


